\section{Discussion and Limitations}
\label{sec:discussion}

\textbf{Workload regimes with limited gains.}
\systemname\ targets the regime in which expert transfer is on (or
near) the critical path. When effective transfer bandwidth is high
enough that the baseline already overlaps most transfers (e.g.,
Qwen2-MoE on H20 in \autoref{sec:evaluation}), the controller
converges to a small~$S$ and contributes only the
$O(L_{\text{MoE}})$ controller overhead, which we measure to be
under $50\,\mu$s per update on a commodity CPU. We did not observe a
configuration in which \systemname\ degraded the baseline, but we do
not claim universal speedups: a model whose entire expert footprint
fits in HBM under the deployment cap, or a workload whose routing is
so concentrated that one or two experts cover it, gains little from
adaptive prefetching.

\textbf{Memory-cap sensitivity.}
Our evaluation fixes a 20\,GB HBM cap to expose expert-residency
contention on every device under study. The cap was chosen to
emulate consumer GPUs and multi-tenant deployments where KV-cache,
activations, and other tenants compete for HBM with expert weights;
it does not represent the only operating point of interest. The
behavior of \systemname\ at smaller caps (10\,GB) is bounded by the
$S_{\min}$ floor of the controller, which prevents thrashing at the
cost of losing some overlap; at larger caps the controller converges
to higher~$S$ and gains diminish as more experts become persistently
resident. Characterizing this curve quantitatively is left to future
work; we expect the qualitative trend to mirror
\autoref{fig:M1}, since the underlying overlap--pollution tradeoff
is the same.

\textbf{Multi-stream and multi-tenant deployment.}
\systemname\ targets a single-stream serving setup with continuous
batching. Co-locating multiple model instances on a single GPU
introduces inter-instance contention on the transfer stream and on
HBM, which the current controller does not explicitly model: each
instance optimizes its own~$S$, and contention manifests in the
shared bandwidth estimate~$C_s$. We expect the per-instance
controller to remain stable (because (C2) bounds per-step change to
$\pm1$), but the joint optimum may not match the union of local
optima. Lifting the controller to a coordinated multi-stream
setting---potentially by sharing $C_s$ and $\tilde{m}_t$ across
instances---is a natural extension.

\textbf{Predictor cold start and online retraining.}
The forecasting layer relies on a CPU-side random forest trained
offline from logged tuples. On a previously unseen model or
substantially shifted workload, the forecaster falls back to the
model's own pregate top-$k$, which we measured matches the
fixed-horizon baseline. The system therefore degrades gracefully
rather than catastrophically. Online retraining (incremental forest
updates as new $(\mathbf{t},\ell,S,\mathbf{a})$ tuples arrive) is
straightforward to add and does not interact with the controller's
stability proof, but we have not characterized its convergence
empirically.

\textbf{Threats to generality.}
Our experiments cover three MoE models from the DeepSeek and Qwen
families and three accelerators spanning a $4\times$ effective-bandwidth
range. We did not evaluate Mixtral-class models (which use a different
expert sizing) or larger DeepSeek-V3-class models; we expect the same
mechanism to apply, but the optimal hyperparameters
($\lambda, \theta_h$) may shift. We also do not evaluate disaggregated
or pipeline-parallel deployments, where prefetch coordination crosses
device boundaries.